\chapter{文献综述}\label{chap:literature-review}

大语言模型智能体（large language model agent，以下简称智能体）已从单轮问答系统发展为能够调用工具、访问外部环境、维护上下文并连续执行任务的复合系统。在这类系统中，基础模型并不单独决定运行结果；提示组织、上下文装配、工具接口、控制循环、状态管理和验证规则也会改变其行为。本文将这组支撑智能体运行的基础设施称为智能体运行框架（agent harness，以下简称运行框架）。开发者在排查问题时，除了确认任务是否成功，还需要说明某个现象如何出现、哪些运行记录可以支持当前判断，以及这些记录与具体运行机制有何关系。因此，本章从系统评价、调试认知、轨迹可观测性（observability）、分布式追踪、分析溯源与人因评价等方面梳理相关研究。

\section{智能体运行框架与运行理解}

\subsection{从模型评价到系统运行评价}

早期大语言模型研究多以准确率、胜率或基准任务得分衡量模型能力。但在真实产品中，用户看到的结果同时受模型、提示、检索、界面、业务规则和协作流程影响。van der Maden等人\cite{ref01}访谈了19名大语言模型产品实践者，其中17人遇到了“结果与可行动性鸿沟”：团队能够收集分数与反馈，却难以据此确定应修改哪个系统环节。Khanna等人\cite{ref04}提出的人工智能行动后复盘方法（After-Action Review for AI，AAR/AI）也显示，将总体表现继续分解到具体决策和内部预测，有助于用户发现已知故障；但其精确率仍然有限。这两项研究表明，系统层面的运行理解不能由最终分数代替，同时也不宜将局部异常直接写成确定根因。

运行框架的相关研究进一步说明了模型评价与系统评价的差别。Chen等人\cite{ref20}将运行框架划分为执行环境与沙箱、工具接口、上下文与记忆、生命周期与编排、可观测、验证评估和治理安全等层次，失败证据往往散布在其中多个位置。Lin等人\cite{ref21}将系统提示、工具、中间件、技能、子智能体、长期记忆和配置视为可编辑的框架组件，并强调这些组件的变更应当可追踪、可比较、可回滚。Amini等人\cite{ref29}在四种运行时上执行相同的智能体声明，仍观察到准确率、延迟和执行行为的差异。所以，本研究所说的“运行理解”，是根据实际执行记录分析运行机制如何发挥作用，而不是仅阅读静态配置或复述模型给出的理由。

\subsection{运行轨迹、机制关联与故障归因}

智能体运行轨迹（execution trace）通常包含模型请求与响应、工具调用与返回值、文件或环境变化、智能体间消息、重试和终止等事件。轨迹记录可以回答“发生了什么”，运行机制则规定系统如何组织这些事件。开发者必须将两者对应起来，才能进一步分析某个现象的成因。Zhang等人\cite{ref03}将多智能体故障归因分为责任智能体与责任步骤两个维度。在127个系统、184个归因任务中，最佳方法识别责任智能体的准确率为53.5\%，定位具体步骤的准确率只有14.2\%。这一结果说明，自动归因可以作为调查线索，但不应不经核查就当作事实。

运行理解与自动修复也有不同的目标。运行框架修复工具HarnessFix\cite{ref20}将原始轨迹与框架实现编译为面向运行框架的轨迹中间表示（Harness-aware Trace Intermediate Representation，HTIR），再通过数据流、控制流和实现锚点将失败步骤映射到责任层次，以便生成和验证修补方案。智能体框架工程方法AHE\cite{ref21}则利用组件可观测、经验可观测和决策可观测形成自动演化闭环。这些工作证明了轨迹与机制之间的关系可以被计算，却没有直接回答开发者如何核对证据、识别推导前提或保留未知结论。本研究因而将人的证据判断与故障定位放在主要位置，不以自动生成修复方案为首要目标。

\subsection{上下文构造作为研究切口}

上下文构造连接模型、工具与历史状态，是运行框架中影响后续行为的关键环节。一次模型请求具体包含哪些系统指令、历史消息、工具结果与检索材料，以及这些内容在组装过程中是否被保留、替换、裁剪或摘要，决定了模型在当前步骤实际能够读取的信息。结构化日志框架AgentTrace\cite{ref26}将上下文事件与认知事件、操作事件并列记录，表明仅保留工具调用无法重建智能体当时的输入条件。系统级可观测工具AgentSight\cite{ref19}同时观测模型侧语义意图与内核侧行为，其结果显示高层意图与低层效果之间存在语义差距。Amini等人\cite{ref29}对开放智能体规范（Open Agent Specification）的跨框架实验又说明，相同声明在不同运行时中仍可能产生不同行为。以上研究为本文提供了一个可操作的切口：追踪信息从何处进入请求、在构造过程中发生了什么变化，以及后续步骤是否仍可读取该信息。

\section{开发者调试与意义建构的认知基础}

\subsection{从异常现象到“为什么”与“为什么没有”}

传统调试研究将开发者的任务区分为发现错误位置与判断、修复错误两个环节。Katz和Anderson\cite{ref15}发现，程序员即使能够修复已指明位置的错误，也未必能根据异常行为顺利找到该位置。这一区分同样适用于智能体：读懂某条工具报错，不等于能够发现导致报错的上下文裁剪或控制分支。

Ko和Myers\cite{ref05,ref06,ref07}围绕“为什么”与“为什么没有”问题（why and why-not questions）开展的一系列研究，将调试入口从猜测代码位置转向针对可见行为提问。其Whyline调试界面允许用户围绕程序输出选择问题，再查看相关运行事件和因果链。早期实验中，该界面缩短了定位时间；后续面向Java程序语言的Whyline研究也提高了参与者完成真实调试任务的比例。这类工具的价值在于，它们以开发者实际遇到的问题组织证据，并保留从问题返回运行记录和代码的路径。

智能体调试工具已经开始采用类似思路。多智能体交互调试工具AGDebugger\cite{ref02}支持暂停、重置和修改历史消息，便于开发者在错误继续传播之前试验其他路径。该研究的14人实验显示，交互调试有助于定位部分问题，但参与者在选择干预位置和修改方式时仍经常失败。软件开发智能体调试工具AgentStepper\cite{ref18}将断点、单步执行、中间状态检查和现场编辑引入智能体轨迹。在12人实验中，参与者识别智能体实现缺陷的成功率从17\%升至60\%，但一般轨迹理解成绩只从64\%升至67\%。由此可见，用户能够操作轨迹，并不代表其已形成正确解释。

\subsection{信息觅食、层级搜索与跨表示协调}

Lawrance等人\cite{ref08}从信息觅食理论（information foraging theory）解释程序员的调试行为。研究发现，开发者会根据错误消息、标识符、文件名和调用关系形成“信息气味”（information scent），并在文件、方法和调用关系构成的信息结构中决定是否继续追踪。Ko等人\cite{ref14}的软件维护研究显示，参与者平均约有35\%的有效时间用于导航、建立关系和收集信息。Rasmussen和Jensen\cite{ref09}对电子设备排故的研究则表明，熟练人员会在系统、功能单元与部件之间切换层级，同时考虑线索收益、测量成本和认知负担。智能体调试界面若仅提供完整日志，开发者仍需要自己完成这些搜索工作；更有用的方式是让用户从异常现象下钻到局部事件，再回溯至相关机制。

视图数量的增加并不会自然改善理解。Romero等人\cite{ref34}在同时呈现代码、输出与图形表示的调试环境中发现，用户的表现还取决于他们能否理解外部表示，并在不同表示之间建立对应关系。Cornelissen等人\cite{ref33}对执行轨迹可视化进行受控实验，在八项程序理解任务中，轨迹可视化组的用时减少22\%，答案正确性提高43\%。该结果说明结构化轨迹可能降低理解成本，但前提是任务、表示与交互方式彼此匹配。在本研究中，时间轴、关系图、实现位置与证据详情需要通过稳定标识和明确的选中状态相互对应，以减少用户在视图之间重复寻找同一事件。

\subsection{数据—框架理论与困惑消解}

对运行过程的理解不是一次性接收事实，而是不断形成和修正解释的过程。Klein等人\cite{ref10,ref11}对意义建构（sensemaking）研究的梳理指出，人会利用既有框架选择和解释信息，新线索又可能强化、质疑或替换原有框架。数据—框架理论（Data–Frame theory）将这种关系描述为数据与框架的双向循环\cite{ref12}。当异常工具结果、缺失消息或意外终止无法纳入当前解释时，开发者需要比较其他可能性，并重新组织证据。

Pirolli和Card\cite{ref17}将意义建构过程分为信息觅食循环和意义建构循环：前者负责搜索、筛选与整理材料，后者负责建立证据结构、产生假设并形成结论。Kadivar等人\cite{ref38}开发的CzSaw可视分析系统记录分析历史、依赖关系与分支，并允许分析者回到旧状态、探索替代路径。Xu等人\cite{ref37}则区分了低层事件、分析动作、子任务和总体任务，指出语义层级越高，系统越难仅依赖日志可靠推断。因而，运行理解界面既要降低查找证据的成本，也要容纳临时假设、反例与尚未解决的问题。

Starr和Storey\cite{ref16}根据27名开发者的访谈，将故障排查概括为困惑、试探、观察和恢复等过程，并指出困惑会占用注意与工作记忆。由此，界面不宜一次堆叠所有日志字段和图形关系，而应围绕当前问题逐步展开相关信息，并允许用户回退、比较和修正判断。

\subsection{以人为中心的可解释性与信任校准}

可解释人工智能（explainable artificial intelligence，XAI）研究区分了“系统能够生成解释”与“用户获得正确理解”。Miller\cite{ref47}结合哲学、认知心理学和社会心理学研究，指出人类解释具有选择性、对比性和社会性，因而特征重要度等单向输出并不必然促进理解。Abdul等人\cite{ref52}对289篇论文及其引用网络的分析也表明，可解释、可问责与可理解系统需要算法研究与人机交互研究配合。对智能体运行界面而言，关键不在于显示尽可能多的内部信息，而在于解释能否对应开发者的问题、返回原始证据并说明推断前提。

国内研究也在从模型透明性转向以人为中心的可理解性。范俊君等人\cite{ref42}认为，智能时代的人机交互需要重新审视心理模型、交互原则与评价框架。范向民等人\cite{ref43}梳理了人机交互与人工智能从交替发展到协同进化的关系。刘伟和孙惟一\cite{ref46}则讨论了人机协同中的确定性与不确定性、事实与价值及多层意义等问题。这些讨论为本研究提供了明确限定：界面不能只展示系统结论，还应让用户识别结论的证据边界。

关于解释内容与交互方式，Amershi等人\cite{ref48}经过多轮专家归纳和验证，提出18条人机交互指南，其中包括明确系统能力边界、支持高效纠错并及时告知行为变化。Liao等人\cite{ref49}用用户可能提出的问题组织可解释需求，发现用户会分别追问为什么发生、为什么没有发生和改变条件后会如何。Wang等人\cite{ref53}将认知与决策理论引入用户中心的可解释框架，强调解释方式应对应用户的决策任务与认知偏差。因此，本研究以开发者的证据查询为交互入口，不预设单一的根因叙事。

解释工具本身也可能产生误导。Kaur等人\cite{ref50}对11名数据科学家的情境调查和197人问卷显示，专业用户也会误用或过度信任可解释工具，并且难以准确说明部分可视化的含义。Ehsan等人\cite{ref51}提出社会透明性（social transparency），将过往决策中的内容、原因、人员和时间信息放回具体组织情境。孔祥维等人\cite{ref41}和李瑶等人\cite{ref44}对人工智能决策可解释性及评价方法的综述同样指出，评价需同时考虑模型属性、应用风险与用户结果，不能只判断系统是否给出了解释。这些研究使“实际记录、规则推导与目前未知”成为值得单独考察的界面变量。

\section{智能体轨迹可观测与交互调试}

\subsection{从结构化日志到多粒度轨迹}

智能体可观测性研究首先需要确定记录范围。AgentTrace\cite{ref26}将日志分为操作、认知和上下文三个平面，并用追踪标识（trace ID）与跨度标识（span ID）组织模型交互、环境动作和上下文变化。该工作为复现、安全分析与问责提供了日志结构，但尚未检验不同呈现方式对开发者理解的影响。AgentSight\cite{ref19}利用扩展伯克利包过滤器（extended Berkeley Packet Filter，eBPF）在不修改智能体代码的情况下捕捉模型通信和系统事件，报告的运行开销低于3\%。但是，内核事件属于直接观测，高层意图及两者的对应关系可能包含规则或模型推断，界面需要让用户看出这种差别。

分布式追踪研究为轨迹组织提供了较成熟的技术基础。X-Trace\cite{ref30}通过显式传播任务标识，将跨应用、协议和管理域的事件重建为任务树。Dapper\cite{ref32}用追踪与跨度表示一次请求的分布式路径，并以公共库插桩和采样控制覆盖范围与运行开销。这些系统说明，可靠的跨组件关联依赖稳定标识和因果上下文传播，而不是事后按时间邻近程度猜测。采样、未插桩组件或标识传播中断都会留下空白，所以“没有记录”不能直接推出“没有发生”。Pivot Tracing\cite{ref31}进一步以先发生关系连接跨组件事件，展示了规则推导的分析价值，也说明推导结果需要保留规则和输入证据。

\subsection{面向人的观察、调试与干预工具}

现有智能体调试工具可按主要交互方式分为三类。第一类以时间和层级结构帮助用户浏览长轨迹。AgentLens\cite{ref22}从原始事件中构建行为结构，用层级时间视图呈现多智能体行为的演变，并支持从整体模式下钻到具体行为及候选原因。AgentGUI\cite{ref13}同时展示任务、轨迹、产物与成本，并提供人工与自动纠偏。在8人实验中，参与者定位关键轨迹信息的平均用时由145秒降至90秒，答案正确率从80\%提高到93\%。这些结果支持以概览与细节协同的方式展示长轨迹；但候选原因和自动管理建议仍然包含推断，应与直接记录区分。

概览需要压缩事件，但压缩可能删去关键参数、异常返回值或消息边界。AgentLens\cite{ref22}使用层级聚合，轨迹图系统Graph of Trace\cite{ref24}用有向图突出阶段与依赖，AgentGUI\cite{ref13}则通过任务和产物视图过滤终端噪声。这些方法能够缩小搜索范围，但被压缩的节点仍需要回到原始事件。若摘要由大语言模型生成，界面还应说明生成范围和可能的遗漏。这与Whyline\cite{ref05,ref07}的设计原则一致：系统可以帮助用户缩小范围，但不能切断结论与原始证据之间的联系。

第二类工具以结构图重建任务与依赖关系。AgentGraph\cite{ref23}将智能体、任务、工具、输入输出与人员组织成知识图谱，每个图元素都可回到对应的轨迹片段。Graph of Trace\cite{ref24}将科学研究智能体的执行事件组织成持续更新的有向图，方便领域专家检查工作流和中间步骤。图结构适合表达长链依赖，但这类研究多以案例展示或主观反馈为主，尚不足以说明用户能否准确区分已记录关系、算法重建关系与未知关系。

第三类工具将观察与干预结合。AutoGen Studio\cite{ref25}用JavaScript对象表示法（JavaScript Object Notation，JSON）描述智能体与工作流，并通过拖拽编辑、运行试验和消息查看支持原型构建与调试。LADYBUG\cite{ref28}支持用户检查每一步的输入输出、在任意步骤干预并仅重算受影响部分，同时利用大语言模型生成候选错误与修正方案。AGDebugger\cite{ref02}和AgentStepper\cite{ref18}进一步提供暂停、回退、单步执行和现场编辑。这类工具降低了比较替代路径的成本，却也将观察、判断和改变系统紧密结合。若用户在原因尚未确认时修改轨迹，新结果可能被误当作对原假设的验证。面向运行理解的界面因而需要先支持证据定位与机制核对，再将干预作为后续操作。

\begin{table}[htbp]
\centering
\caption{代表性智能体运行可观测与交互调试系统比较}
\label{tab:report-1}
\scriptsize
\renewcommand{\arraystretch}{1.18}
\begin{tabularx}{\textwidth}{>{\raggedright\arraybackslash}p{2.65cm} *{5}{>{\raggedright\arraybackslash}X}}
\toprule
\textbf{代表工作} & \textbf{主要表示} & \textbf{核心交互} & \textbf{轨迹与机制关联} & \textbf{证据状态呈现} & \textbf{人因评价特点} \\
\midrule
AGDebugger\cite{ref02} & 多智能体消息轨迹 & 暂停、重置、编辑消息 & 主要关联消息与智能体角色 & 未显式区分 & 形成性访谈与14人实验，关注定位和干预 \\
AgentGUI\cite{ref13} & 任务、轨迹、产物与成本视图 & 观察、人工与自动纠偏 & 可定位轨迹元素，机制映射有限 & 未系统区分 & 8人实验，测量时间、正确率和工作负担 \\
AgentStepper\cite{ref18} & 结构化会话与代码变化 & 断点、单步执行、编辑、检查状态 & 关联智能体动作与实现执行 & 未显式区分 & 12人对照研究，测量轨迹理解、缺陷识别和负担 \\
AgentLens\cite{ref22} & 层级时间视图 & 概览、筛选、下钻、原因追踪 & 部分原因由算法挖掘 & 推断边界不突出 & 使用情景与用户研究 \\
AgentGraph\cite{ref23} & 交互式知识图谱 & 图浏览、返回原始轨迹 & 图元素链接精确轨迹片段 & 可核验，状态分类有限 & 以系统展示为主 \\
Graph of Trace\cite{ref24} & 实时有向工作流图 & 过程检查与阶段回顾 & 依赖由监测与解析模块重建 & 未独立编码 & 领域专家主观评价 \\
HarnessFix\cite{ref20} & 面向运行框架的中间表示 & 自动归因与修复 & 数据流、控制流与实现锚点 & 面向机器诊断 & 以基准性能和修复效果为主 \\
\bottomrule
\end{tabularx}
\end{table}

\subsection{自动归因、修复与人的理解}

智能体故障诊断已从轨迹展示进入轨迹解释阶段。Who\&When\cite{ref03}尝试自动确定责任智能体与关键步骤，但步骤级准确率仍然较低。HarnessFix\cite{ref20}将轨迹片段映射到运行框架层次与实现对象，并在四类基准上将初始框架表现提高6.3至18.4个百分点。AHE\cite{ref21}为每次运行框架编辑记录预期结果，并用下一轮任务表现对其验证。它们的共同贡献是将失败轨迹处理为可定位机制问题的结构化证据，而不只是一个负向结果。

然而，修复后表现改善，不足以证明开发者已理解当初的故障。Parnin和Orso\cite{ref35}对自动故障定位工具的研究指出，算法将错误语句排在前列，不代表程序员能够理解根因；真实调试还需要沿数据和控制依赖导航、检查状态并通过重运行验证假设。智能体归因界面若只给出“责任步骤”，却不展示数据来源、关联规则与缺失信息，用户可能更快地采取操作，也可能更快地得出过度确定的结论。

人机界面与自动修复系统对输出的要求也不相同。自动修复系统可以在大量候选方案中搜索，再用回归测试筛选补丁；开发者界面则需要说明当前判断依据什么，以及判断不成立时应在何处继续调查。AHE\cite{ref21}将修改前的预测与下一轮结果绑定，HarnessFix\cite{ref20}利用实现锚点约束修复范围。对人机交互界面而言，相应的做法是让每项故障判断都成为可验证的主张：注明支持它的记录与规则，并保留尚未排除的替代解释。

\section{轨迹与机制关联的证据基础}

\subsection{溯源信息的对象与用途}

溯源信息（provenance）用于描述数据或结论从何处而来、经历了什么过程，以及为何形成当前结果。Ragan等人\cite{ref36}将可视化和数据分析中的溯源信息分为数据、可视化、交互、洞见与理由等类型，并强调记录对象应与使用目的相匹配。复现运行可能只需保存参数、输入和执行环境；若要说明开发者为何得出某项判断，则还需要记录假设、注释和证据选择过程。Wang等人\cite{ref27}针对大语言模型智能体归纳了执行溯源框架，将检索证据、工具结果、记忆项、环境观察、中间声明、动作与最终答案组织为有类型的关系，同时指出统一模式、语义级溯源与恢复导向评价仍然不足。

分析溯源研究还揭示了语义层级带来的推断风险。Xu等人\cite{ref37}从低层事件、分析动作、子任务到总体任务建立层次，认为高层任务与意图难以仅靠系统日志准确推断。CzSaw\cite{ref38}通过自动记录交互历史并允许用户编辑、分支和重放，在机器记录与人的解释之间建立了分工。对运行框架而言，工具返回文本、请求长度和消息删除属于可直接记录的事实；“某条结果因上下文裁剪而未进入后续请求”需要结合数据流规则推导；“模型为何忽略仍然可见的信息”则可能超出现有观测能力。界面若将三者以相同方式呈现，就会掩盖其语义层级和证据强度的差异。

\subsection{实际记录、规则推导与目前未知}

结合分布式追踪、智能体可观测性与溯源研究，本文将运行理解中的证据关系分为三类状态。“实际记录”指系统在运行时直接捕获，且能够返回原始事件的事实，例如某次模型请求包含的消息列表、工具调用参数和文件写入事件。“规则推导”指依据明示规则连接多个已记录事实，例如根据追踪标识重建父子跨度，根据先发生关系确定依赖，或根据上下文构造函数将输入字段映射到请求片段\cite{ref30,ref31,ref32}。“目前未知”指必需信号因未插桩、采样、未持久化或超出模型可观测边界而缺失，现有证据不足以支持肯定结论。

这三种状态不等于简单的高、中、低可信度。直接记录可能不完整，时间戳也可能有误；规则推导可以十分可靠，但需向用户公开规则与前提；在必需信号缺失时，“未知”是对证据边界的准确描述。Sacha等人\cite{ref39}指出，不确定性会从数据、模型和可视化继续传播到人的解释。Hullman\cite{ref40}调查了90名可视化作者并访谈13名设计者，发现作者虽承认不确定性的价值，却常因担心界面复杂、说服力下降和读者负担而省略这类信息。因此，证据状态应当显示，但呈现方式不宜用密集警告干扰主要任务。

“目前未知”也不应被折叠成视觉上的次要信息。当界面呈现流畅的候选原因或醒目的高亮路径时，用户可能提前停止寻找反证。数据—框架理论认为，无法被当前框架解释的线索是触发框架重构的重要条件\cite{ref12}。因而，证据状态设计至少需要让用户区分事实与推导，在推导处展开规则、输入与反例条件，并在未知处说明缺少哪类信号及可能的补证方式。

\begin{table}[htbp]
\centering
\caption{智能体运行框架理解中的证据状态}
\label{tab:report-2}
\scriptsize
\renewcommand{\arraystretch}{1.18}
\begin{tabularx}{\textwidth}{>{\raggedright\arraybackslash}p{2.0cm} *{4}{>{\raggedright\arraybackslash}X}}
\toprule
\textbf{证据状态} & \textbf{判定依据} & \textbf{典型例子} & \textbf{界面需回答的问题} & \textbf{主要风险} \\
\midrule
实际记录 & 有原始事件、时间戳与稳定标识，可返回源记录 & 请求中实际包含的消息；工具真实返回值；文件写入 & 记录位于何处，是否完整，属于哪次运行 & 将“记录到结果”误解为“已证明原因” \\
规则推导 & 由已记录事实和公开规则计算，能说明输入与规则 & 由父子跨度还原调用链；由代码锚点推导上下文来源 & 使用了什么规则，前提是否满足，能否复算 & 将算法推断误当作直接观察 \\
目前未知 & 必需信号缺失或超出系统可观测边界 & 未保存的中间上下文；模型内部为何选择某动作 & 缺少什么，是否能通过补充插桩确认 & 用时间邻近、相关性或语言流畅性填补空白 \\
\bottomrule
\end{tabularx}
\end{table}

\subsection{从可追溯关系到可辩护结论}

轨迹与机制关联的作用不在于为关系图增加更多连线，而在于让用户能够沿连线检查证据。X-Trace\cite{ref30}和Dapper\cite{ref32}说明，稳定的跨组件关联需要在执行过程中传播上下文。Pivot Tracing\cite{ref31}表明，事前未记录的指标可以通过动态插桩和关系查询补充，但查询结果仍依赖明确的因果顺序。AgentGraph\cite{ref23}让图元素返回精确轨迹片段，HarnessFix\cite{ref20}则将运行步骤连接到运行框架实现对象与修复算子。综合这些工作，本研究关注三类关联：表示信息来源、传递与变换的数据流，表示分支、重试、校验和终止的控制流，以及连接配置项、函数或文件的实现锚点。

这三类关联在证据性质上并不相同。界面需要标明某条连线表示直接观测到的数据传递、按规则重建的控制依赖，还是运行事件与实现位置的静态映射。例如，日志中出现工具结果，只能说明该结果确实生成。若后续请求快照表明该结果未进入模型输入，才能继续判断上下文传递发生了中断。若请求快照缺失，则应保留未知，不宜仅根据后续输出反推该结果已被丢弃。

\section{现有研究的评价方法}

\subsection{从系统展示到受控理解任务}

智能体可观测性研究常见的评价方式有系统展示、专家评议和受控实验。AutoGen Studio\cite{ref25}、AgentGraph\cite{ref23}、LADYBUG\cite{ref28}和AgentTrace\cite{ref26}主要用系统展示或使用情景证明工具可以接入真实框架并呈现复杂轨迹。这类评价可以检验技术可行性，但不能说明用户是否因此形成了更准确的机制解释。Graph of Trace\cite{ref24}通过领域专家反馈评价可理解性与可用性，AgentLens\cite{ref22}结合使用情景与用户研究考察行为分析。专家访谈适合发现工作流与界面问题，但主观清晰感可能与客观正确率不一致。

具有对照条件和客观答案的实验更适合检验运行理解。Whyline系列研究\cite{ref05,ref06,ref07}使用预设故障与可判断的任务完成标准比较调试表现。Cornelissen等人\cite{ref33}比较了常规集成开发环境与增加轨迹可视化后的完成时间和正确性。AAR/AI\cite{ref04}预先确认了十个故障，再分别测量用户识别结果的召回率与精确率。AgentGUI\cite{ref13}和AgentStepper\cite{ref18}则分别测量轨迹信息定位、缺陷发现、时间与工作负担。本研究需要沿用这种评价思路，直接比较参与者的结论与已知运行机制，而不只使用可用性量表。

\subsection{指标选择与实验条件隔离}

结合上述研究，运行理解的评价可以包括四类结果。理解正确性用于衡量参与者对事件顺序、数据来源、机制作用和故障位置的判断。证据定位用于衡量其是否找到关键轨迹片段、运行框架实现对象及二者关系。结论校准包括无依据断言、将推导误述为事实，以及在证据不足时保留判断。过程与负担指标则包括完成时间、导航路径、视图切换、任务负担指数（NASA Task Load Index，NASA-TLX）与主观信任。若只报告系统性能，会遗漏用户为理解结果支付的时间和认知成本；若只报告满意度，又无法证明理解是否正确。

实验设计还需要分开信息数量与信息组织方式的影响。Parnin和Orso\cite{ref35}指出，自动调试研究常把算法排名质量直接当作对程序员的帮助。Romero等人\cite{ref34}的研究说明，增加表示可能同时增加信息量与协调负担。Sacha等人\cite{ref39}和Hullman\cite{ref40}的工作则提示，不确定性标记既可能帮助校准信任，也可能增加认知成本。如果实验组同时获得更多底层日志、更强的关联算法和更丰富的界面，就无法判断改善来自哪一部分。为此，各条件应使用同一数据底座，依次比较常规轨迹、增加轨迹与机制关联的界面，以及在相同关联上增加证据状态的界面。

\subsection{客观真值、任务材料与生态效度}

运行理解实验的难点之一是建立可核对的标准答案。真实项目具有较高的生态效度，却常缺少唯一、可验证的根因；人工小程序易于精确控制，又可能将调试简化为寻找明显错误。AAR/AI\cite{ref04}在同一系统中预设十个已知故障，并同时测量召回与误报。AgentStepper\cite{ref18}结合真实软件开发智能体与事先分析的轨迹缺陷，在保留任务复杂度的同时为缺陷定位提供判断依据。Cornelissen等人\cite{ref33}使用代表性软件系统和八项程序理解任务，避免以单一问题代表整体理解。对本研究而言，较合适的材料组合是真实运行记录与经过仪器化处理的构造案例。

真实运行可用于检验界面在噪声、长轨迹和不完整数据中的表现；构造案例则可有计划地操纵上下文保留、替换、裁剪、摘要和工具结果注入等机制，使关键问题具有明确答案。任务应覆盖事件顺序、信息来源、变换机制和证据充分性等层次，并在难度、轨迹长度、关键证据位置与机制类型上进行配平。除最终答案外，研究还应记录参与者打开了哪些事件、是否查看机制依据、在何处修改判断，以及证据不足时是否保留未知。

实验也需要控制学习效应与界面熟悉度。如果同一参与者依次体验三种条件，后续条件可能因对任务更熟悉而受益；如果采用组间设计，小样本又容易受经验差异影响。可采用拉丁方平衡条件顺序，并在正式任务前安排不计分的训练案例。Romero等人\cite{ref34}发现，外部表示知识会显著影响多视图调试表现。因此，研究还需记录智能体开发经验、传统调试经验和可视化使用经验，再结合屏幕记录、交互日志、任务后访谈与客观答案，区分界面未提供线索、参与者未发现线索和发现线索后仍推断错误三种情况。

\begin{table}[htbp]
\centering
\caption{智能体运行理解研究的评价范式}
\label{tab:report-3}
\scriptsize
\renewcommand{\arraystretch}{1.18}
\begin{tabularx}{\textwidth}{>{\raggedright\arraybackslash}p{2.0cm} *{4}{>{\raggedright\arraybackslash}X}}
\toprule
\textbf{评价范式} & \textbf{代表研究} & \textbf{常用结果} & \textbf{可回答的问题} & \textbf{主要局限} \\
\midrule
系统展示或案例 & AutoGen Studio\cite{ref25}、AgentGraph\cite{ref23}、LADYBUG\cite{ref28} & 功能覆盖、案例完成、性能开销 & 系统是否可实现并接入真实轨迹 & 不能证明人的理解得到改善 \\
专家评议或访谈 & Graph of Trace\cite{ref24}、AgentLens\cite{ref22} & 可用性、可理解性、需求反馈 & 界面是否符合工作实践 & 主观清晰感可能偏离客观正确性 \\
受控调试实验 & Whyline\cite{ref05,ref06,ref07}、AgentStepper\cite{ref18} & 成功率、时间、缺陷定位、任务负担 & 交互工具是否改善具体调试任务 & 样本和任务通常较小，生态效度有限 \\
真值驱动的理解评价 & AAR/AI\cite{ref04}、轨迹可视化实验\cite{ref33} & 召回率、精确率、正确率、用时 & 用户结论是否与已知事实一致 & 需要精心构造且难度相当的任务材料 \\
\bottomrule
\end{tabularx}
\end{table}

\section{国内外研究现状与不足}

国外相关研究已经形成从传统程序调试、分布式追踪和可视分析到智能体交互工具的技术脉络。Whyline\cite{ref05}、信息觅食\cite{ref08}和故障定位\cite{ref15}研究提出了以用户问题组织证据、用导航行为解释调试成本的方法。X-Trace\cite{ref30}、Dapper\cite{ref32}和Pivot Tracing\cite{ref31}提供了跨组件关联、上下文传播与动态因果查询的技术。溯源信息\cite{ref36,ref37}与不确定性研究\cite{ref39,ref40}进一步将记录历史、语义层级、信任校准和证据边界纳入界面设计。近年的AGDebugger\cite{ref02}、AgentStepper\cite{ref18}、AgentGUI\cite{ref13}、AgentLens\cite{ref22}和AutoGen Studio\cite{ref25}开始针对大语言模型智能体长轨迹提供观察、下钻与干预能力。

国内研究在人机交互理论、人工智能决策可解释性和人机协同评价方面已经形成一定基础\cite{ref41,ref42,ref43,ref44,ref46}，但直接面向智能体长轨迹与运行框架调试的实证工作仍较少。近期相关技术研究更多关注轨迹结构化、运行框架诊断、自动演化与科学工作流可视化。Who\&When\cite{ref03}用大规模多智能体任务数据揭示了自动故障归因的难度。HarnessFix\cite{ref20}利用数据流、控制流与实现锚点将失败证据连接到运行机制，AHE\cite{ref21}则从组件、经验和决策三个层面建立可观测闭环。AgentLens\cite{ref22}、AgentGraph\cite{ref23}与Graph of Trace\cite{ref24}分别探索了层级时间视图、知识图谱和实时工作流图对复杂智能体行为的表达；Wang等人\cite{ref27}试图用统一的溯源视角组织检索、工具、记忆与故障证据。现有工作已经具备从技术可观测性转向交互理解的基础，但针对开发者认知过程的受控人因证据仍然较少。

对上述研究的比较显示了四个与本课题直接相关的不足。现有可观测工具主要按线性日志、时间视图或任务图组织运行事件，轨迹现象与运行框架机制之间的显式映射仍不充分；即使HarnessFix\cite{ref20}已能生成这种映射，其主要使用者也是自动诊断与修复模块，而不是开发者。原始记录、规则推导与模型生成的候选解释常被放在相近的视觉层级，用户难以判断结论的证据性质与未知边界。已有评价多集中于任务成功率、系统性能、可用性或主观清晰度，很少在同一研究中测量机制理解正确性、证据定位和无依据结论。此外，部分研究在实验条件中同时增加日志、关联算法与界面元素，不利于区分信息量和信息组织方式各自的作用。

基于这些不足，本研究考察一次智能体运行已经发生、暂不要求系统自动修复时，开发者如何通过交互界面将轨迹现象与运行框架机制建立关联，并识别该关联属于实际记录、规则推导还是目前未知。研究将先通过形成性研究分析开发者的问题、搜索路径与误判，再在现有运行架构原型上增加轨迹与机制关联及证据状态，最后利用具有客观真值的真实运行和仪器化构造案例开展三条件受控实验。三种条件使用同一底层数据，依次呈现常规轨迹、轨迹与机制关联，以及附带证据状态的关联，以分别检验关联组织和证据状态的作用。

\section{本章小结}

现有研究已经为智能体运行理解提供了轨迹记录、跨组件关联、层级浏览和交互干预等方法。调试与意义建构研究表明，开发者需要从可观察异常出发，逐步搜索证据并修正解释\cite{ref05,ref08,ref12,ref17}。分布式追踪与溯源研究则说明，稳定标识、因果上下文与可核验历史是建立关联的基础，但记录缺失不能作为事件未发生的证据\cite{ref30,ref31,ref36,ref37}。智能体可观测与交互调试工具已能记录、浏览、下钻和干预长轨迹，但轨迹事件、运行机制和证据性质仍未被统一呈现\cite{ref02,ref18,ref23}。

本研究因而把问题收敛到上下文构造过程，考察开发者如何将运行现象与数据流、控制流和实现锚点建立关联，并判断关联属于实际记录、规则推导还是目前未知。后续研究将以理解正确性、证据定位、无依据结论、完成时间与任务负担为主要指标，检验不同界面呈现方式的作用。
